\newcommand{\figuraEnlacesSuperficie}{%
\begin{figure}[H]
\centering
\begin{tikzpicture}[>=Latex,font=\small,line cap=round,line join=round]
  \fill[blue!8] (0.8,0.5) rectangle (11.2,4.2);
  \draw[very thick,blue!65!black] (0.8,4.2)--(11.2,4.2);
  \node at (9.9,4.75) {vacío};
  \foreach \x in {1.6,2.8,...,10.0}
    \foreach \y in {1.0,2.0,3.0,4.0}
      \fill[orange!75] (\x,\y) circle (0.14);
  \foreach \x in {1.6,2.8,...,8.8}
    \foreach \y in {1.0,2.0,3.0,4.0}
      \draw[orange!45!black] (\x+0.14,\y)--(\x+1.06,\y);
  \foreach \x in {1.6,2.8,...,10.0}
    \foreach \y in {1.0,2.0,3.0}
      \draw[orange!45!black] (\x,\y+0.14)--(\x,\y+0.86);
  \draw[very thick,red!75!black,dashed] (4.0,4.15)--(4.0,5.15)
    node[above] {enlace faltante};
  \draw[very thick,red!75!black] (4.0,4.0) circle (0.27);
  \draw[very thick,green!45!black] (8.8,2.0) circle (0.27);
  \node[red!75!black,anchor=west] at (4.35,3.72) {molécula superficial};
  \node[green!45!black,anchor=west] at (9.15,2.0) {molécula interior};
  \node[align=center] at (4.0,0.05) {superficie:\\[1mm]cinco enlaces};
  \node[align=center] at (8.8,0.05) {interior:\\[1mm]seis enlaces};
\end{tikzpicture}
\caption{Modelo de red para la energía superficial. Una molécula superficial
carece de un enlace respecto de una molécula interior, por lo que posee una
energía positiva adicional por unidad de área.}
\label{fig:enlaces-superficie}
\end{figure}%
}